\section{Conclusion and Future Work}

\subsection{Achievements}

This project set out to answer a precise question: can a faithful miniature
NFVI --- one that genuinely meets the demands of a 5G core --- be built on a
single 16~GB laptop, and can a real 5G core run on it? The answer is yes. The
concrete achievements are:

\begin{itemize}
    \item An \textbf{NFVI platform} (three-node k3s cluster) provisioned
          end-to-end as Infrastructure-as-Code, acting as the NFVI and VIM of
          the ETSI model.
    \item A \textbf{multi-plane network} that gives the UPF dedicated N3/N6
          user-plane interfaces (Multus + macvlan + Whereabouts), strictly
          isolated from management and cluster-signalling traffic.
    \item A \textbf{real 5G Standalone core} (Open5GS) deployed as CNFs, with a
          simulated RAN and subscriber (UERANSIM), demonstrated end to end:
          NF registration, UE registration, PDU session, and \emph{actual
          user-plane traffic through the UPF}.
    \item \textbf{Elastic scaling} of a stateless NF and a RAM-optimised
          \textbf{observability} stack giving live visibility into the core.
    \item A report whose every design decision is traced back to a concrete
          requirement of the 5G core (Table~\ref{tab:requirements}).
\end{itemize}

The key point is the framing: this is not ``three nodes running Kubernetes''.
It is an infrastructure platform whose every capability exists because a 5G
network function needs it --- and that claim is validated by running the
network function itself.

\subsection{Limitations}

The testbed is deliberately scoped to one host, with honest consequences.
The control plane runs on a \textbf{single k3s server node}, a SPOF (Sec.~4.3).
The UPF uses the \textbf{kernel GTP-U data path}, so throughput reflects the
virtualised environment rather than the line rate that SR-IOV/DPDK would give.
MongoDB uses \texttt{emptyDir}, so subscriber data is \textbf{not persistent}.
The RAN is \textbf{simulated} (UERANSIM), not real radio hardware, and
orchestration (NFVO) is manual (Ansible/Helm) rather than a full MANO stack.

\subsection{Future Work}

\begin{itemize}
    \item \textbf{Highly-available control plane} --- add two more k3s server
          nodes (embedded etcd) to remove the node0 SPOF.
    \item \textbf{Accelerated data plane} --- give the UPF SR-IOV or DPDK
          interfaces for line-rate, production-like packet processing.
    \item \textbf{Persistent subscriber store} --- back MongoDB with a
          persistent volume and a real provisioning workflow.
    \item \textbf{GitOps orchestration} --- replace manual Helm with Argo CD /
          Flux to approximate the NFVO role declaratively.
    \item \textbf{Network slicing and multiple PDU sessions} --- exercise
          several S-NSSAIs and DNNs to model realistic 5G services.
\end{itemize}

In summary, the project demonstrates that the essential structure and
behaviour of a cloud-native 5G core --- and the NFV infrastructure it
requires --- can be reproduced faithfully on commodity hardware, providing an
accessible, end-to-end platform for learning and experimentation in 5G and NFV.
